\begin{abstract}
Multi-stage contact manipulation requires controlling object motion, orientation and tool contact beyond intermediate position attainment. We investigate constrained parameterized residual reinforcement learning in a 16-environment JAKA testbed using Isaac Sim 5.1 and Isaac Lab 2.3.0. Historical ordinary residual PPO improves aggregate unseen-geometry completion over feedback, but most improvement comes from the foundational pushing task rather than guided-slot insertion (P2). After preserving that evidence and a negative effect-history pilot, three bounded physical research iterations test contact-target authority, shared rotated-contact geometry with tool yaw, and a stateful re-contact executor. The last iteration compares ordinary and execution-active Gaussian PPO under the same 87-dimensional observation, five-dimensional action interface, controller, reward and interaction budget. A saved-mask objective demonstrably changes actor score gradients on identical real rollout samples. However, three independently trained policies obtain 22/25/20 versus 24/24/18 accepted P2 completions out of the same 32 development cases; shared strong feedback obtains 26/32. The paired gains of +2/-1/-2 do not support a consistent method advantage. The implementation isolates a real actor-update mechanism, not a new marginalization theorem or validated recovery policy. Actual paired videos, process traces and independently audited outcomes expose both physical successes and failures. Development scenes are not a sealed final test, and state-based synthetic manipulation is not field or real-robot validation. A publishable novel-method claim remains unsupported.
\end{abstract}
\noindent\textbf{Keywords:} contact manipulation; residual reinforcement learning; stateful skill parameters; execution-aware policy objectives; reproducible robotics.
\vspace{4pt}
\noindent\fbox{\parbox{0.94\columnwidth}{\small \textbf{Evidence status.} True online PhysX training and actor-gradient isolation are verified. The new three-seed P2 objective does not consistently improve task completion or outperform shared strong feedback. No sealed-final-test, visual-policy or hardware-safety claim is made.}}
